\documentclass[11pt]{article}
\usepackage[margin=1in]{geometry}
\usepackage{amsmath,amssymb,booktabs,xcolor,tcolorbox}
\tcbuselibrary{skins}
\newtcolorbox{idea}{colback=blue!4,colframe=blue!40,title=In plain words}
\newtcolorbox{quickcheck}{colback=green!4,colframe=green!40!black,title=Quick check}
\newcommand{\1}{\mathbf{1}}
\title{Simplified proof 1: the covariance law\\[4pt]
\large every step written out}
\author{Folder 22 / simplified\_proofs --- companion to \texttt{proofs/1\_exact\_statements.pdf}, Proposition 3}
\date{}
\begin{document}\maketitle

\begin{idea}
We want to know: if we increase one layer-2 weight a little (say $c_2$), does the loss go down?
The answer we will prove is: \textbf{yes, exactly when that weight moves attention towards keys that
hold the correct next token more often than the keys attention is already on.} The rest of this
document derives that sentence, one small step at a time.
\end{idea}

\section{What we are looking at}
Fix \textbf{one query position} $i$. To keep the notation light, we drop the index $i$ wherever we can.
\begin{itemize}
\item There are $n$ earlier positions (``keys''), numbered $j=1,\dots,n$.
\item Key $j$ holds the token $x_j$ (one of the 8 tokens).
\item The correct next token for our query is $y$.
\item $\mathrm{hit}_j=1$ if key $j$ holds the correct token ($x_j=y$), and $0$ otherwise.
\item Each key gets a \textbf{score} $s_j$ from layer 2. We will only change one weight, called $\kappa$
(it can be $a_2$, $b_2$ or $c_2$).
\end{itemize}

\section{Step 1: scores become attention (softmax)}
The attention on key $j$ is
\[
A_j=\frac{e^{s_j}}{\sum_{k=1}^n e^{s_k}} .
\]
Two facts we will use:
\begin{itemize}
\item every $A_j$ is between 0 and 1, and
\item they add up to one: $\sum_j A_j=\frac{\sum_je^{s_j}}{\sum_ke^{s_k}}=1$. So the $A_j$ behave like probabilities.
\end{itemize}

\section{Step 2: how attention changes when one score changes}
We need $\frac{\partial A_k}{\partial s_j}$: how the attention on key $k$ moves if we nudge the score of key
$j$. Write $E=\sum_m e^{s_m}$, so $A_k=e^{s_k}/E$, and note that $\frac{\partial E}{\partial s_j}=e^{s_j}$.

\paragraph{Case $k=j$.} Quotient rule on $e^{s_j}/E$:
\[
\frac{\partial A_j}{\partial s_j}=\frac{e^{s_j}\cdot E-e^{s_j}\cdot e^{s_j}}{E^2}
=\frac{e^{s_j}}{E}-\Big(\frac{e^{s_j}}{E}\Big)^2=A_j-A_j^2=A_j(1-A_j).
\]
\paragraph{Case $k\neq j$.} Now the numerator $e^{s_k}$ does not depend on $s_j$; only $E$ does:
\[
\frac{\partial A_k}{\partial s_j}=-\frac{e^{s_k}\cdot e^{s_j}}{E^2}=-A_kA_j .
\]
\paragraph{Both cases in one formula.} Using $\1[k=j]$ (1 if $k=j$, else 0):
\[
\boxed{\ \frac{\partial A_k}{\partial s_j}=A_k\big(\1[k=j]-A_j\big)\ }
\]
(check: for $k=j$ this gives $A_j(1-A_j)$, for $k\ne j$ it gives $-A_kA_j$). Raising one key's score
takes attention away from all the other keys, in proportion to how much attention they had.

\section{Step 3: the prediction is a mix of the key tokens}
Layer 2 copies tokens (that is what $W_{ov}=I$ means). So the predicted distribution is
\[
f=\sum_{k=1}^nA_k\,e_{x_k},
\]
where $e_{x_k}$ is the one-hot vector of token $x_k$. Its entry for the correct token $y$ is
\[
f[y]=\sum_kA_k\,\1[x_k=y]=\sum_kA_k\,\mathrm{hit}_k .
\]
\begin{idea}
$f[y]$ is the attention-weighted fraction of keys that hold the right token: \textbf{the current
average correctness}.
\end{idea}

\section{Step 4: how the correct-token probability changes when one score changes}
Differentiate $f[y]=\sum_kA_k\mathrm{hit}_k$ with respect to $s_j$, using Step 2 ($\mathrm{hit}_k$ does not depend
on scores):
\[
\frac{\partial f[y]}{\partial s_j}=\sum_k\mathrm{hit}_k\,A_k\big(\1[k=j]-A_j\big)
=\underbrace{\mathrm{hit}_j\,A_j}_{\text{the }k=j\text{ term}}\;-\;A_j\underbrace{\sum_kA_k\mathrm{hit}_k}_{=f[y]} .
\]
So
\[
\boxed{\ \frac{\partial f[y]}{\partial s_j}=A_j\,\big(\mathrm{hit}_j-f[y]\big)\ }
\]
\begin{idea}
Raising key $j$'s score helps the correct token if key $j$ is \textbf{better than average}
($\mathrm{hit}_j=1>f[y]$) and hurts if it is \textbf{worse than average} ($\mathrm{hit}_j=0<f[y]$).
\end{idea}

\begin{tcolorbox}[colback=orange!5,colframe=orange!60!black,title=What ``better than average'' means --- a worked example]
\textbf{The ``average''} is $f[y]$: of all the attention, the share that sits on keys holding the right
token. A key is \textbf{better than average} if it is a hit (it is 100\% right, which beats any mix that
is less than 100\% right), and \textbf{worse than average} if it is a miss (0\% right).

\medskip
\textbf{Why raising a score works like this.} Raising $s_j$ does two things at once (Step 2):
\begin{enumerate}
\item key $j$ \textbf{gains} attention;
\item that attention is \textbf{taken from all the other keys}, in proportion to what they had, so it
comes from a pool that is about $f[y]$ right.
\end{enumerate}
So attention moves \emph{from a mix that is $f[y]$ right} \emph{to key $j$}. If key $j$ is a hit,
the right share goes up. If it is a miss, it goes down.

\medskip
\textbf{Numbers.} Four keys, attention $A=(0.4,\,0.4,\,0.1,\,0.1)$, hits $=(1,\,0,\,1,\,0)$, so
$f[y]=0.4+0.1=0.5$ (half the attention is on right keys).

\smallskip
\emph{Raise key 1 (a hit) by $0.1$.} Key 1 gains $0.4\cdot0.6\cdot0.1=0.024$. The others lose
$A_k\cdot0.4\cdot0.1$: key 2 loses $0.016$, key 3 $0.004$, key 4 $0.004$. The change in
$f[y]$ counts only the hits: $+0.024$ (key 1) $-0.004$ (key 3) $=\mathbf{+0.020}$.
Formula: $A_1(\mathrm{hit}_1-f[y])\cdot0.1=0.4\cdot(1-0.5)\cdot0.1=+0.020$. \checkmark

\smallskip
\emph{Raise key 2 (a miss) by $0.1$.} Key 2 gains $0.024$, but it is a miss, so that adds nothing. The hits
lose: key 1 $0.016$, key 3 $0.004$. Change in $f[y]$: $\mathbf{-0.020}$.
Formula: $0.4\cdot(0-0.5)\cdot0.1=-0.020$. \checkmark

\medskip
\textbf{Two edge cases.} If \emph{every} key is a hit ($f[y]=1$), raising any score changes nothing
($\mathrm{hit}_j-f[y]=0$): you can't do better than all right. A key with little attention
($A_j$ small) moves $f[y]$ only a little: that is the factor $A_j$ in front.
\end{tcolorbox}

\section{Step 5: from the prediction to the loss}
\paragraph{The readout.} The model's final probabilities are
\[
q=\frac{f+b}{Z},\qquad Z=\sum_{t=0}^{7}\big(f[t]+b[t]\big),
\]
where $b$ is the small fixed bias (it gives every token a little probability).

\paragraph{$Z$ is a constant.} Since $f$ is a mix of one-hot vectors with weights adding to 1, its entries
add to 1: $\sum_tf[t]=\sum_kA_k\sum_t\1[x_k=t]=\sum_kA_k\cdot1=1$. Therefore $Z=1+\sum_tb[t]$, which does
not depend on the scores at all.

\paragraph{The loss.} $\ell=-\log q[y]=-\log\big(f[y]+b[y]\big)+\log Z$. Only the first term changes with
the scores, and $\frac{d}{du}(-\log u)=-\frac1u$, so
\[
\frac{\partial\ell}{\partial f[y]}=-\frac{1}{f[y]+b[y]}=:-w .
\]
We name $w=\frac1{f[y]+b[y]}$ the \textbf{surprise}: it is large when the model gives the correct token
little probability.

\paragraph{Chain rule (loss $\leftarrow$ $f[y]$ $\leftarrow$ score $s_j$).} Combine with Step 4:
\[
\boxed{\ \frac{\partial\ell}{\partial s_j}=\frac{\partial\ell}{\partial f[y]}\cdot\frac{\partial f[y]}{\partial s_j}
=-w\,A_j\,\big(\mathrm{hit}_j-f[y]\big)\ }
\]

\section{Step 6: from a score to the weight $\kappa$}
Layer 2's score is
\[
s_j=a_2\,\nu_j+b_2\,\sigma_i\,r_j[x_i]+c_2\,\langle r_i,r_j\rangle .
\]
Each weight appears \textbf{once}, multiplied by something that does not contain that weight. So the
derivative of the score with respect to a weight is just its multiplier, which we call the
\textbf{feature} $\phi_j$:
\begin{center}
\begin{tabular}{lll}\toprule
weight $\kappa$ & feature $\phi_j=\partial s_j/\partial\kappa$ & meaning\\\midrule
$a_2$ & $\nu_j$ & key $j$ is a noise token\\
$b_2$ & $\sigma_i\,r_j[x_i]$ & key $j$ came after my token\\
$c_2$ & $\langle r_i,r_j\rangle$ & key $j$ has the same recovered state as me\\\bottomrule
\end{tabular}
\end{center}
\paragraph{Why nothing else moves.} These three weights live in layer 2. Layer 1 does not use them, so
$r_i,r_j$ do not change when $\kappa$ changes. The \emph{only} thing $\kappa$ changes is the scores $s_j$.

\paragraph{Chain rule again (loss $\leftarrow$ all scores $\leftarrow$ $\kappa$).} Changing $\kappa$ changes every score,
so we add up over all keys:
\[
\frac{\partial\ell}{\partial\kappa}=\sum_j\frac{\partial\ell}{\partial s_j}\cdot\frac{\partial s_j}{\partial\kappa}
=\sum_j\Big(-w\,A_j\,(\mathrm{hit}_j-f[y])\Big)\phi_j .
\]
Multiply by $-1$ (we want the direction that \emph{lowers} the loss):
\[
\boxed{\ -\frac{\partial\ell}{\partial\kappa}=w\sum_jA_j\,\big(\mathrm{hit}_j-f[y]\big)\,\phi_j\ }
\]

\section{Step 7: why this sum is a covariance}
Treat the attention $A_j$ as probabilities: pick a random key $J$ with $\Pr(J=j)=A_j$. Then
\begin{itemize}
\item the average correctness is $\mathbb E[\mathrm{hit}_J]=\sum_jA_j\mathrm{hit}_j=f[y]$ (Step 3);
\item the covariance of two quantities is $\mathrm{Cov}(X,Y)=\mathbb E\big[(X-\mathbb EX)(Y-\mathbb EY)\big]$.
\end{itemize}
Our sum is $\mathbb E\big[(\mathrm{hit}_J-f[y])\,\phi_J\big]$. It equals the covariance, because subtracting the mean of
$\phi$ changes nothing:
\[
\mathbb E\big[(\mathrm{hit}_J-f[y])(\phi_J-\mathbb E\phi)\big]
=\mathbb E\big[(\mathrm{hit}_J-f[y])\phi_J\big]-\mathbb E\phi\cdot\underbrace{\mathbb E\big[\mathrm{hit}_J-f[y]\big]}_{=f[y]-f[y]=0}.
\]
So, finally,
\[
\boxed{\ -\frac{\partial\ell}{\partial\kappa}=w\cdot\mathrm{Cov}_{J\sim A}\big(\mathrm{hit}_J,\ \phi_J\big)\ }
\]
and for the whole loss (the average over query positions $i$) we just average:
$-\frac{\partial L}{\partial\kappa}=\frac1N\sum_iw_i\,\mathrm{Cov}_i(\mathrm{hit},\phi)$.

\begin{idea}
\textbf{Covariance $>0$}: the keys the weight lifts are more often correct than the current average,
so gradient descent \textbf{increases} the weight. \textbf{Covariance $<0$}: it lifts mostly wrong keys,
so gradient descent \textbf{decreases} it. The surprise $w\ge0$ never changes the sign; it only makes
positions where the model is wrong count more.
\end{idea}

\section{A worked example}
One query, 4 keys, $b[y]=0.01$:
\begin{center}
\begin{tabular}{ccccc}\toprule
key $j$ & $A_j$ & $\mathrm{hit}_j$ & $\phi_j$ & $A_j(\mathrm{hit}_j-f[y])\phi_j$\\\midrule
1 & 0.4 & 1 & 1 & $0.4\times0.5\times1=0.20$\\
2 & 0.4 & 0 & 0 & $0$\\
3 & 0.1 & 1 & 1 & $0.1\times0.5\times1=0.05$\\
4 & 0.1 & 0 & 0 & $0$\\\bottomrule
\end{tabular}
\end{center}
\begin{enumerate}
\item $f[y]=0.4+0.1=0.5$ (Step 3).
\item $w=1/(0.5+0.01)\approx1.96$.
\item Sum of the last column $=0.25$, so $-\partial\ell/\partial\kappa\approx1.96\times0.25=0.49>0$: \textbf{raise} $\kappa$.
\end{enumerate}
Change the feature to $\phi=(0,1,0,1)$ (it now lifts the two wrong keys): the column becomes
$0,\ 0.4\times(-0.5),\ 0,\ 0.1\times(-0.5)$, sum $-0.25$, so \textbf{lower} $\kappa$.

\begin{quickcheck}
(1) If every key has $\phi_j=1$, what is the covariance? (Hint: subtract the mean of $\phi$.) Why does that
make sense --- what does raising $\kappa$ do to the attention if every score goes up by the same amount?\\
(2) Where in the derivation did we use that layer 2 \emph{copies} tokens? (Hint: Step 5.)
\end{quickcheck}

\section{Why this matters for the rest of the project}
For the state weight $c_2$, the feature is ``key $j$ has the same recovered state as me''.
\begin{itemize}
\item \textbf{After a noise token}, my recovered state \emph{is} the current state, so the lifted keys came after
the current state: they are its successors, often correct. Covariance $>0$, so $c_2$ goes \textbf{up}.
\item \textbf{After a signal token}, my recovered state is the \emph{previous} state, so the lifted keys are
successors of the previous state: mostly wrong. Covariance $<0$, so $c_2$ goes \textbf{down}.
\end{itemize}
That is the tug-of-war of figure \texttt{f4}: one law, two opposite signs.
\end{document}
